% folio.tex — the anatomy of a folio. Copy to src/<subject>/<name>.tex and replace
% every part. One subject, one A4 page, dense facts + pictures, no noise.
%
% Header comment (keep it — the build and the reader both use it):
%   line 1      what this folio covers, in one line
%   Sources:    where every fact can be checked (docs, specs, papers, man pages)
%   @labels:    key=value pairs, e.g. area=<subject> kind=<what> platform=<if any>;
%               discussion=<n> once the folio's thread exists in GitHub Discussions
%               (hidden on the page; "Discuss this folio" opens thread n)
%   @tags:      comma list; reuse tags other folios already use
%
% Sources: <url or reference>, <url or reference>
%   (the fact-check trail, for whoever edits the folio — readers get "Further reading")
% @labels: area=c-systems kind=memory
% @tags: stack, heap, malloc, lifetime, dangling-pointer
\documentclass[8pt]{extarticle}
\usepackage{printup-sheet}

% Linked text: the website links it, paper prints the text only (paper holds no links).
% Defined here, not in the style: the website's renderer reads macros from the folio itself.
\newcommand\link[2]{\pdfonly{#1}\htmlonly{\href{#2}{#1}}}

\begin{document}

\sheettitle{Stack vs heap}{c-systems · folio}

% One sentence that would be worth keeping if it were the only thing kept.
\oneliner{The \textbf{stack} holds a function's locals and dies with its frame
(LIFO, automatic, fast); the \textbf{heap} holds what \texttt{malloc} returns and
lives until \texttt{free} (manual, slower, fragmentable).}

% A picture first when the subject has structure: flows, layers, comparisons.
\begin{multicols}{2}

\section{How it works}
\begin{itemize}
  \item Each call pushes a \emph{frame}: return address, saved registers, locals.
        Return = move the stack pointer back — nothing is erased.
  \item The stack is small (8~MiB main thread on Linux, often 4--16~KiB on an
        RTOS task) and grows \emph{down}; overflow is a crash, not an error code.
  \item \texttt{malloc} asks the allocator for a block; the allocator asks the
        OS (\texttt{brk}/\texttt{mmap}) only when its free lists run dry —
        \link{glibc}{https://man7.org/linux/man-pages/man3/malloc.3.html} gives a
        block over 128~kB (the default threshold) its own \texttt{mmap}.
\end{itemize}

\section{Example}
\begin{lstlisting}[language=C]
int *bad(void)  { int x = 42; return &x; }   // dangling: frame gone
int *good(void) { int *p = malloc(sizeof *p);
                  if (p) *p = 42; return p; } // lives until free(p)
\end{lstlisting}

\columnbreak

\section{Picture}
\begin{tikzpicture}[sheet]
  \node[font=\bfseries\small] at (0,3.1) {stack};
  \node[cell, minimum width=22mm, fill=sheetBlue!10] (m) at (0,2.5) {main()};
  \node[cell, minimum width=22mm, fill=sheetBlue!18] (g) at (0,1.95) {good(): p};
  \node[cell, minimum width=22mm, draw=sheetRed, dashed] (b) at (0,1.4) {bad(): x (gone)};
  \draw[flow] (-1.45,2.7) -- node[left, note]{grows} (-1.45,1.2);
  \node[font=\bfseries\small] at (3.3,3.1) {heap};
  \node[cell, minimum width=16mm, fill=sheetGreen!15] (h) at (3.3,1.95) {42};
  \draw[hot] (g.east) -- (h.west);
\end{tikzpicture}

% Misconceptions stated as facts — not interview questions.
\section{Traps}
\begin{itemize}
  \trap{``The stack is faster memory'' — no: same RAM. Allocation is faster
        (one pointer bump) and it stays hot in cache.}
  \trap{Returning \texttt{\&local} often \emph{works} in a test — the frame
        is not wiped until the next call reuses it.}
\end{itemize}

\section{Remember}
\textbf{Stack = scope, heap = \texttt{free}.} Whoever calls \texttt{malloc} owns the \texttt{free}.

% Web only, last before Related: a few links worth reading, one line each on why.
% Paper prints nothing here and keeps the space.
\begin{htmlonly}
\section{Further reading}
\begin{itemize}
  \item \href{https://man7.org/linux/man-pages/man3/malloc.3.html}{malloc(3)} — the 128~kB
        \texttt{mmap} threshold, and the extra arenas glibc adds when threads contend
\end{itemize}
\end{htmlonly}

\end{multicols}

\noindent{\footnotesize\color{sheetGrey}\textit{Related:} memory leaks · \texttt{alloca} / VLAs · RAII (C++) · ownership (Rust)}

\end{document}
